\documentclass[10pt]{article}
\usepackage[margin=0.9in]{geometry}
\usepackage{amsmath,microtype}
\title{Guide: The Verified Residue-Coverage Tool}
\author{Collatz proof project}
\date{September 5, 2026}
\setlength{\emergencystretch}{3em}
\begin{document}
\raggedright
\maketitle
\section*{What has been proved}
Every positive Collatz orbit visits a number whose remainder on division
by nine is two. Such visits occur after any specified time. This is a
formalization of a known theorem of Monks and coauthors, not a new proof
of the Collatz conjecture.

\section*{How the argument works}
First move past any multiples of three. The remaining remainders are
one, two, four, five, seven, and eight. If the orbit avoids two, the
allowed remainder transitions force it eventually to stay at eight.
An even number with remainder eight moves to remainder four, so this
would require only odd steps forever. Our existing Lean parity argument
rules that out for natural starting values.

\section*{How to use it in Lean}
Import \texttt{Collatz.ResidueCoverage}.
The theorem \texttt{exists\_two\_mod\_nine} supplies a hitting time.
The theorem \texttt{arbitrarily\_late\_two\_mod\_nine} also accepts
a lower bound on the time.
The theorem \texttt{reaches\_one\_of\_two\_mod\_nine} reduces full
convergence to convergence for positive seeds in that progression.
Its convergence hypothesis still needs a proof.

\section*{The missing connection}
Earlier work proves that every hypothetical unbounded orbit contains
starting points whose multiplicative coefficients never fall below one.
It also restricts the smallest such starting point.
The residue-coverage theorem does not say its chosen visit has that
coefficient property. Two separately guaranteed kinds of visits need
not occur at the same time. A proof using both results must establish
that extra connection or use a different invariant.

\section*{Verification}
From \texttt{lean/}, run
\texttt{lake build Collatz.ResidueCoverage}.
From the repository root, run
\texttt{python3 analysis/audit\_theorems.py --build} with
the project's Lean executable on \texttt{PATH}.
The selected audit checks the actual theorem statements and axiom
dependencies; no additional axiom asserts the literature result.
\end{document}
